Zur Überprüfung der Zielerreichung und zur Beantwortung der formulierten Fragestellungen wird eine quantitative Evaluierung der entwickelten Ansätze durchgeführt. 

Grundlage der Evaluation bildet ein repräsentativer Referenzdatensatz (Ground Truth), der strukturierte Zielwerte für ausgewählte Produktdaten enthält. Auf dieser Basis erfolgt 
ein systematischer Vergleich beider Ansätze hinsichtlich der Qualität der erzeugten Daten sowie der Robustheit gegenüber variierenden Eingabeformaten.

Die Auswertung erfolgt anhand definierter Key Performance Indicators (KPIs), die sowohl die Qualität als auch die Strukturtreue abbilden. Dabei orientiert sich die Auswahl an 
etablierten Dimensionen der Datenqualität \parencite[vgl.][S. 5--33]{wang1996beyond}:

\begin{itemize}
    \item \textit{Schema Compliance Rate:} Anteil der generierten Datensätze, die dem vorgegebenen Datenbankschema entsprechen (Dimension der 
    \textit{Representational Consistency}).
    \item \textit{Vollständigkeit (Coverage):} Anteil der korrekt befüllten Datenfelder im Verhältnis zu allen erforderlichen Feldern 
    \parencite[Dimension der \textit{Completeness}, vgl.][S. 24]{batini2016data}.
    \item \textit{Konsistenz:} Grad der strukturellen und semantischen Einheitlichkeit der erzeugten Daten. Diese Metrik validiert die Widerspruchsfreiheit innerhalb der 
    extrahierten Datensätze \parencite[vgl.][S. 31]{batini2016data}.
    \item \textit{Exact Match Rate:} Strikte Übereinstimmung der generierten Werte mit dem Ground Truth. Diese binäre Metrik wird häufig zur Evaluation von Extraktionsaufgaben 
    in strukturierten Benchmarks herangezogen \parencite[vgl.][S. 1283]{data_extraction_ai}.
    \item \textit{Precision, Recall und F1-Score:} Aggregierte Maße der Extraktionsgüte, wobei Precision die Genauigkeit und Recall die Trefferquote beschreibt 
    \parencite[vgl.][S. 67f.]{data_extraction_ai}.
    \item \textit{Effizienzmetriken:} Ergänzend werden Effizienzaspekte wie Laufzeit und Verarbeitungskosten betrachtet, was insbesondere im betrieblichen Kontext von hoher 
    Relevanz ist \parencite[vgl.][S. 3]{ferrara2014web}.
\end{itemize}

Abschließend werden die Ergebnisse interpretiert und hinsichtlich ihrer praktischen Einsetzbarkeit sowie ihrer Übertragbarkeit auf vergleichbare Anwendungsfälle bewertet.